\sgaStableTarget{sga6.e06.prop.6-2}\textbf{Corollary 6.2.} \uline{With the hypotheses of} 6.1, \uline{let} @@M0@@ \uline{be a sub}-@@M1@@-\uline{ring generated by the classes of finitely many sheaves}. \uline{Then the} @@M2@@-\uline{ring filtration of} @@M3@@ \uline{is discrete}, \uline{and} @@M4@@ \uline{is nilpotent}. \uline{Moreover}, @@M5@@ \uline{is a module of finite type over} @@M6@@ (V 3.9.1).

Since @@M7@@ is quasi-compact, there is a finite covering of @@M8@@ by disjoint open sets such that, on each open set of this covering, the sheaves whose classes generate @@M9@@ all have constant rank. The standard argument then reduces us to the case where these sheaves have constant rank on @@M10@@; @@M11@@ is then augmented to @@M12@@.

Let @@M13@@, @@M14@@, be the sheaves whose classes generate @@M15@@, @@M16@@ their classes, and @@M17@@. We can write @@M18@@ as a quotient of the @@M19@@-ring generated by the @@M20@@ subject to the relations @@M21@@ for @@M22@@, that is, @@M23@@; consequently, @@M24@@ is noetherian. But by Proposition 6.1, the @@M25@@, @@M26@@ are nilpotent. Thus the augmentation ideal of @@M27@@ is nilpotent; moreover, it is precisely @@M28@@.

Every element of @@M29@@ can be written as a polynomial with no components of weight @@M30@@ in the @@M31@@, @@M32@@. Let @@M33@@ be the largest integer such that there exists @@M34@@ such that @@M35@@; let @@M36@@ be an integer such that the @@M37@@-th powers of the @@M38@@, @@M39@@ are all zero. For @@M40@@, one sees at once that @@M41@@; the filtration of @@M42@@ is therefore discrete.

Since the filtration of @@M43@@ is discrete, there are only finitely many nonzero monomials in the @@M44@@, @@M45@@. These monomials form a system of generators of @@M46@@ as a @@M47@@-module, which completes the proof.

